\section{The auxiliary field}\label{sec:auxiliary-field}

The top weights have indices in a cyclic group $H$. Summing over $H$ in
the prime sum isolates $G$ up to three quadratic terms. The quadratic
fields stay fixed on a power window. The order of $H$ depends on the
point $x$ in that window.

\subsection{The field and its primes}

Let $q,d$ be distinct primes with
\[
  q\equiv3\pmod 8,\qquad d\equiv7\pmod{16},\qquad d\equiv-1\pmod q.
\]
For $\ell=q,d$ write $\chi_{-\ell}$ for the primitive quadratic
character of $\mathbb Q(\sqrt{-\ell})$. Its value at $a$ is the
Legendre symbol $(a/\ell)$. For $e\ge1$ put $n=2^e$, $M=4n$ and
$\zeta_M=\exp(2\pi i/M)$. Set
\[
  K_e=\mathbb Q(\zeta_M+\zeta_M^{-1}),\qquad
  E=K_e(\sqrt{-d},\sqrt{-q}),\qquad
  H=(\mathbb Z/M\mathbb Z)^\times/\{\pm1\}.
\]

\begin{proposition}\label{prop:aux-field}
The group $H$ is cyclic of order $n$. The field $E$ is a CM field with
maximal real subfield $K_e(\sqrt{dq})$ and Galois group
\[
  \mathcal G=(\mathbb Z/2\mathbb Z)^2\times H.
\]
In coordinates $(i,j,h)$ the actions on $\sqrt{-d}$ and $\sqrt{-q}$
have signs $(-1)^i$ and $(-1)^j$, respectively. Complex conjugation
is $c=(1,1,0)$.
\end{proposition}

\begin{proof}
The units congruent to $1$ modulo $4$ represent $H$ bijectively.
There are $n$ such units, and $5$ generates them, since
$v_2(5^{2^r}-1)=r+2$ for $r\ge0$, by successive squaring.
In $\mathbb Q(\zeta_{Mqd})$ the Chinese remainder theorem gives
independent actions at $M,d,q$. On the first factor the kernel of
restriction to $K_e$ is $\{\pm1\}$. On each odd factor the kernel
of the Legendre symbol is the subgroup of squares. This gives
$\mathcal G$ and the two stated signs. Conjugation fixes $K_e$ and
negates both square roots. Its fixed field $K_e(\sqrt{dq})$ is totally
real and has index two in $E$.
\end{proof}

Fix an embedding of $E$ in $\mathbb C$ and identify its embeddings with
$\mathcal G$. Exactly one of $g,c+g$ fixes $\sqrt{-d}$. The embeddings
that fix it therefore form a CM type
\[
  \Phi=\{(0,j,h)\},\qquad f=|\Phi|=2n.
\]
We write the group law additively. For $t\in H$, the translation
\[
  p_t(1,j,h)=(0,j+1,h+t)
\]
matches the complementary type with $\Phi$. In
Proposition~\ref{prop:accounting} the nodes run over this $\Phi$,
$v=\tfrac12|\operatorname{disc}(E)|^{1/(4n)}$, and
$\lambda_1=L'(1,\chi_{-d})/L(1,\chi_{-d})$.

For each character $\psi$ of $\mathcal G$ let $\chi_\psi$ denote the
primitive Dirichlet character that induces $\psi$ on
$(\mathbb Z/Mqd\mathbb Z)^\times$. Let $I_p$ be the inertia subgroup
at $p$ and let $\sigma_p$ lift arithmetic Frobenius in
$\mathcal G/I_p$. For $k\ge1$ character orthogonality gives
\begin{equation}\label{eq:field-frobenius}
  \frac{1}{4n}\sum_{\psi\in\widehat{\mathcal G}}
       \overline{\psi(g)}\,\chi_\psi(p^k)
  =\frac{\mathbf 1_{\{g\in k\sigma_p+I_p\}}}{|I_p|}.
\end{equation}
Indeed, averaging $\psi$ on $k\sigma_p+I_p$ gives $\psi(k\sigma_p)$
if $\psi$ is trivial on inertia and zero otherwise. These values are
exactly $\chi_\psi(p^k)$, including zero when $p$ divides its conductor.
Fourier inversion gives \eqref{eq:field-frobenius}.

For $p\nmid Mqd$ the Frobenius is $(i,j,[p])$, where
$(-1)^i=\chi_{-d}(p)$ and $(-1)^j=\chi_{-q}(p)$.
Thus $k\sigma_p$ lies outside $\Phi$ precisely when $p$ is inert in
$\mathbb Q(\sqrt{-d})$ and $k$ is odd. Both $2$ and $q$ split in that
quadratic field, since $-d\equiv1$ modulo $8$ and modulo $q$.
Their Frobenius cosets lie in $\Phi$ for every power. At $d$, only
$\mathbb Q(\sqrt{-d})$ ramifies among the three defining factors of
$E$, so $I_d=\{0,(1,0,0)\}$. Each coset $k\sigma_d+I_d$ has one
element outside $\Phi$, of mass $1/2$ in \eqref{eq:field-frobenius}.

\subsection{The four characters}

For each nonprincipal primitive character $\eta$ put
\[
  G_\eta(y)=\sum_{r\le y}\Lambda(r)\eta(r)(y/r-1)
               +y\,\frac{L'(1,\eta)}{L(1,\eta)},\qquad y\ge1.
\]
Set $G_1=G$ and take the inverse Fourier transform
\[
  F_y(g)=\frac{1}{4n}\sum_{\psi\in\widehat{\mathcal G}}
                   \overline{\psi(g)}\,G_{\chi_\psi}(y).
\]
For the fibre $(1,0,H)$ this gives
\begin{equation}\label{eq:field-projection}
  4\sum_{h\in H} F_y(1,0,h)
   =G(y)+G_{\chi_{-q}}(y)-G_{\chi_{-d}}(y)
      -G_{\chi_{-q}\chi_{-d}}(y).
\end{equation}
Indeed, the sum over $h$ vanishes unless $\psi$ is trivial on $H$.
The four remaining characters are $\chi_{-d}^a\chi_{-q}^b$ with
$a,b\in\{0,1\}$, where exponent zero omits the factor. Their values at
$(1,0,0)$ are $(-1)^a$. The sum over $H$ contributes $n$ in each case.

\subsection{Choice on a power window}

\begin{proposition}\label{prop:field-window}
Fix $0<\varepsilon<1/52$ and $J>1$. For every sufficiently large
$U$, one can choose $q,d$ as above so that
\begin{equation}\label{eq:field-window-bound}
  |G_\eta(y)|\le y^{1/2+\varepsilon}\quad(U\le y\le U^J),\qquad
  \eta\in\{\chi_{-q},\chi_{-d},\chi_{-q}\chi_{-d}\}.
\end{equation}
For each $x\in[U,U^J]$ let $n$ be the least power of two at least
$x^{2+8\varepsilon}$, and form $E$ with these same primes. Then
\begin{equation}\label{eq:field-large-arguments}
  \max_{\psi\in\widehat{\mathcal G}}\frac{|G_{\chi_\psi}(y)|}{y}
    \le \exp(-x^{\varepsilon/8})
    \qquad\bigl(y\ge\tfrac13\exp(x^{\varepsilon/3})\bigr).
\end{equation}
\end{proposition}

\begin{proof}
Put $z=U^{\varepsilon/5}$ and $z'=z\log^2 U$. Apply the two
prime-selection theorems \citep[Theorems~5.10 and~5.12]{adriazola-grh}
with $N_0=8$, $K=J$ and $\tau=\varepsilon$. The first selects $q$,
and the second selects $d$ for that $q$, with
\[
  z/2<\log q\le z,\qquad z'/2<\log d\le z',\qquad
  d\equiv-1\pmod p\quad\text{for every odd prime }p\le z.
\]
Thus all primes $p\le z$ split in $\mathbb Q(\sqrt{-d})$.
Taking the trivial character modulo one as the initial twist gives
\eqref{eq:field-window-bound} and $|\lambda_1|\ll J\log U$. The same
theorems exclude real zeros in $(1-1/400,1)$ for each of the three
characters twisted by a real primitive character of conductor dividing
$8$.

Every real character of $H$ factors through reduction modulo $8$,
since $[5]$ generates $H$ and its image modulo $8$ has order two.
Thus the zero exclusion applies to every real $\chi_\psi$ whose
conductor is divisible by $q$ or $d$. The other real nonprincipal
characters have conductor dividing $8$. Their nonvanishing at $1$
gives a fixed $c_*>0$ such that all the real nonprincipal characters
under consideration have no real zero in $(1-c_*,1)$.

Write $t=\log y$ and $L=\log(Mqd)$. Each conductor divides $Mqd$.
The large-argument estimate \citep[Proposition~5.3 and the following
paragraph]{adriazola-grh}, including its principal-character case,
now gives, with an absolute $c>0$,
\[
  \frac{|G_{\chi_\psi}(y)|}{y}
  \ll L^5\exp\!\bigl(-c\min\{t/L,\sqrt t\}\bigr)
  \qquad(y\ge2).
\]
Since $U\le x$ and $n<2x^{2+8\varepsilon}$, we have
$L\ll x^{\varepsilon/5}\log^2 x$. In the range of
\eqref{eq:field-large-arguments}, $t\ge\tfrac12 x^{\varepsilon/3}$ for
large $U$. Consequently
\[
  t/L\gg x^{2\varepsilon/15}/\log^2 x,\qquad
  \sqrt t\gg x^{\varepsilon/6}.
\]
Both exponents $2/15$ and $1/6$ exceed $1/8$, so the exponential
dominates $L^5$ and gives \eqref{eq:field-large-arguments}.
\end{proof}

At the scale $x$, \eqref{eq:field-projection} and
\eqref{eq:field-window-bound} give
\[
  \left|4\sum_{h\in H}F_x(1,0,h)-G(x)\right|
     \le 3x^{1/2+\varepsilon}.
\]
The weights must recover this sum over $h$ from the excess of one shift.
